\documentclass[11pt,a4paper]{article}
\usepackage[utf8]{inputenc}
\usepackage[T1]{fontenc}
\usepackage{geometry}
\geometry{margin=1in}
\usepackage{hyperref}
\usepackage{enumitem}
\usepackage{booktabs}
\usepackage{longtable}
\usepackage{xcolor}
\usepackage{titlesec}
\usepackage{amsmath,amssymb}

\titleformat{\section}{\Large\bfseries\color{blue!70!black}}{}{0em}{}
\titleformat{\subsection}{\large\bfseries\color{blue!50!black}}{}{0em}{}
\titleformat{\subsubsection}{\normalsize\bfseries\color{blue!30!black}}{}{0em}{}

\title{\textbf{Replication Plan}\\[0.5em]
\large Updated Virophage Taxonomy and Distinction from Polinton-like Viruses\\[0.3em]
\normalsize OSTI ID: 2475938}
\author{Auto-generated Replication Blueprint}
\date{\today}

\begin{document}
\maketitle
\tableofcontents
\newpage

%---------------------------------------------------------------------
\section{Paper Overview}
%---------------------------------------------------------------------
This paper provides an updated taxonomy of virophages---small viruses that parasitize giant viruses---and establishes criteria to distinguish them from related Polinton-like viruses (PLVs). The authors collect virophage and PLV genome sequences from public databases, identify conserved marker genes using HMM profiles, perform phylogenetic analyses (particularly of the major capsid protein (MCP) and packaging ATPase), analyze gene content and genome organization, and propose quantitative taxonomic thresholds (e.g., amino acid identity cutoffs) for genus/family-level classification.

%---------------------------------------------------------------------
\section{Detailed Workflow}
%---------------------------------------------------------------------

\subsection{Phase 1: Genome Collection}
\begin{enumerate}[label=\textbf{1.\arabic*}]
  \item \textbf{Download known virophage genomes} from NCBI GenBank.
    \begin{itemize}
      \item Sputnik, Mavirus, Zamilon, and all others listed in the paper's supplementary tables.
    \end{itemize}
  \item \textbf{Download Polinton-like virus genomes.}
  \item \textbf{Download giant virus genomes} (Mimiviridae) for context.
  \item \textbf{Retrieve metagenomic virophage sequences} from IMG/VR or other databases.
\end{enumerate}

\subsection{Phase 2: Gene Identification}
\begin{enumerate}[label=\textbf{2.\arabic*}]
  \item \textbf{Predict ORFs} in all genomes (Prodigal or GeneMarkS).
  \item \textbf{Identify marker genes} using HMM profiles:
    \begin{itemize}
      \item Major capsid protein (MCP).
      \item Packaging ATPase.
      \item Minor capsid protein, protease, integrase.
    \end{itemize}
  \item \textbf{Run HMMER} (\texttt{hmmsearch}) against custom HMM profiles.
  \item \textbf{Annotate all ORFs} using BLAST against NR or InterProScan.
\end{enumerate}

\subsection{Phase 3: Phylogenetic Analysis}
\begin{enumerate}[label=\textbf{3.\arabic*}]
  \item \textbf{Align marker protein sequences} using MAFFT or MUSCLE.
  \item \textbf{Trim alignments} with trimAl or Gblocks.
  \item \textbf{Construct phylogenetic trees} using:
    \begin{itemize}
      \item Maximum likelihood: IQ-TREE or RAxML with model selection.
      \item Bootstrap support ($\geq$ 1000 replicates) or SH-aLRT.
    \end{itemize}
  \item \textbf{Visualize trees} with iTOL or FigTree.
  \item \textbf{Root trees} using appropriate outgroups.
\end{enumerate}

\subsection{Phase 4: Genome Comparison and Taxonomy}
\begin{enumerate}[label=\textbf{4.\arabic*}]
  \item \textbf{Compute pairwise amino acid identity (AAI)} between all genomes.
  \item \textbf{Compute gene content similarity} (shared orthologous groups).
  \item \textbf{Apply taxonomic thresholds} proposed in the paper.
  \item \textbf{Generate genome synteny plots} to compare gene order.
  \item \textbf{Distinguish virophages from PLVs} based on the established criteria.
\end{enumerate}

\subsection{Phase 5: Validation}
\begin{enumerate}[label=\textbf{5.\arabic*}]
  \item Reproduce phylogenetic trees and verify topology.
  \item Reproduce AAI distribution plots and taxonomic assignments.
  \item Reproduce gene content comparison tables.
\end{enumerate}

%---------------------------------------------------------------------
\section{Datasets}
%---------------------------------------------------------------------

\begin{longtable}{p{4.5cm} p{3.5cm} p{6cm}}
\toprule
\textbf{Dataset / Source} & \textbf{Type} & \textbf{Location / Notes} \\
\midrule
\endfirsthead
\toprule
\textbf{Dataset / Source} & \textbf{Type} & \textbf{Location / Notes} \\
\midrule
\endhead
NCBI GenBank virophage genomes & Genome sequences & \url{https://www.ncbi.nlm.nih.gov/genbank/} \newline Accession numbers from paper's supplementary \\
\midrule
IMG/VR database & Metagenomic viral genomes & \url{https://img.jgi.doe.gov/vr/} \\
\midrule
Pfam / InterPro HMM profiles & Protein family models & \url{https://www.ebi.ac.uk/interpro/} \\
\midrule
Custom virophage HMM profiles & From paper & Supplementary materials or author correspondence \\
\midrule
Paper supplementary tables & Genome lists, accessions & Accompanying the paper \\
\bottomrule
\end{longtable}

%---------------------------------------------------------------------
\section{Models, Tools, and Software}
%---------------------------------------------------------------------

\begin{longtable}{p{4.5cm} p{2.5cm} p{6.5cm}}
\toprule
\textbf{Tool / Library} & \textbf{Version} & \textbf{Purpose / URL} \\
\midrule
\endfirsthead
\toprule
\textbf{Tool / Library} & \textbf{Version} & \textbf{Purpose / URL} \\
\midrule
\endhead
Prodigal & $\geq 2.6$ & Gene prediction \newline \url{https://github.com/hyattpd/Prodigal} \\
\midrule
HMMER & $\geq 3.3$ & HMM-based protein search \newline \url{http://hmmer.org/} \\
\midrule
MAFFT & $\geq 7.4$ & Multiple sequence alignment \newline \url{https://mafft.cbrc.jp/alignment/software/} \\
\midrule
trimAl & $\geq 1.4$ & Alignment trimming \newline \url{http://trimal.cgenomics.org/} \\
\midrule
IQ-TREE & $\geq 2.1$ & Maximum likelihood phylogenetics \newline \url{http://www.iqtree.org/} \\
\midrule
RAxML-NG & latest & Alternative ML tree builder \newline \url{https://github.com/amkozlov/raxml-ng} \\
\midrule
BLAST+ & $\geq 2.12$ & Sequence similarity search \newline \url{https://blast.ncbi.nlm.nih.gov/} \\
\midrule
ETE3 & latest & Tree manipulation and visualization (Python) \newline \url{http://etetoolkit.org/} \\
\midrule
iTOL & web & Interactive tree visualization \newline \url{https://itol.embl.de/} \\
\midrule
CompareM / pyani & latest & AAI and ANI computation \newline \url{https://github.com/dparks1134/CompareM} \\
\midrule
Python (BioPython) & $\geq 3.8$ & Sequence manipulation \newline \url{https://biopython.org/} \\
\bottomrule
\end{longtable}

\subsection{Hardware Requirements}
\begin{itemize}
  \item Standard workstation; phylogenetic analyses can be parallelized.
  \item $\geq$ 16\,GB RAM; $\geq$ 50\,GB storage for genomes and databases.
\end{itemize}

\end{document}
